\section{4.10. Flujo completo}\label{flujo-completo}

\subsection{4.10. Flujo completo del
análisis}\label{flujo-completo-del-anuxe1lisis}

Una vez descritos los distintos módulos que componen EVMAudit, resulta
posible analizar el funcionamiento global de la solución.

El proceso completo implementado por la herramienta se compone de varias
fases consecutivas que permiten transformar un contrato inteligente en
un informe consolidado de vulnerabilidades.

\subsubsection{4.10.1. Recepción del
contrato}\label{recepciuxf3n-del-contrato}

El proceso comienza con la recepción del contrato inteligente que se
desea analizar.

El contrato puede proceder de distintos orígenes:

\begin{itemize}
\tightlist
\item
  ejecución directa de la librería
\item
  aplicación web desarrollada
\item
  futuras integraciones externas.
\end{itemize}

Tras recibir el contrato, EVMAudit detecta automáticamente el nombre del
contrato y configura la versión adecuada del compilador Solidity.

\subsubsection{4.10.2. Obtención de resultados mediante Slither y
Mythril}\label{obtenciuxf3n-de-resultados-mediante-slither-y-mythril}

La siguiente fase consiste en ejecutar las herramientas de análisis
principales.

En primer lugar se ejecuta Slither, aprovechando las ventajas del
análisis estático.

Posteriormente se ejecuta Mythril, incorporando las capacidades
derivadas de la ejecución simbólica.

La combinación de ambas técnicas permite obtener una cobertura superior
a la proporcionada por cada herramienta de forma individual.

!TODO: referencia cruzada a la Sección 2.X (herramientas de análisis).

\subsubsection{4.10.3. Normalización de
resultados}\label{normalizaciuxf3n-de-resultados}

Las salidas generadas por las herramientas se transforman posteriormente
a una estructura común.

Este proceso elimina las diferencias existentes entre ambas herramientas
y permite trabajar con una representación homogénea de las
vulnerabilidades.

La normalización constituye el paso previo necesario para aplicar el
mecanismo de correlación.

\subsubsection{4.10.4. Correlación de
vulnerabilidades}\label{correlaciuxf3n-de-vulnerabilidades}

Una vez normalizados los resultados, EVMAudit agrupa aquellas
vulnerabilidades equivalentes detectadas por varias herramientas.

La correlación permite:

\begin{itemize}
\tightlist
\item
  reducir duplicidades
\item
  aumentar la confianza de determinados hallazgos
\item
  simplificar la interpretación del informe final.
\end{itemize}

Esta fase representa uno de los principales elementos diferenciadores de
la solución desarrollada.

\subsubsection{4.10.5. Generación de propiedades y
fuzzing}\label{generaciuxf3n-de-propiedades-y-fuzzing}

A partir de las vulnerabilidades correlacionadas, el sistema consulta el
catálogo SWC y genera automáticamente las propiedades necesarias para
ejecutar Echidna.

Posteriormente se construye un wrapper Solidity específico y se realiza
una fase adicional de validación mediante fuzzing.

Este proceso proporciona información complementaria sobre determinadas
vulnerabilidades detectadas durante las fases anteriores.

\subsubsection{4.10.6. Generación del informe
final}\label{generaciuxf3n-del-informe-final}

Finalmente, toda la información obtenida durante el análisis es
consolidada en un informe único.

El informe incluye:

\begin{itemize}
\tightlist
\item
  vulnerabilidades detectadas
\item
  herramientas que las han identificado
\item
  nivel de confianza asociado
\item
  resultados del fuzzing
\item
  metadatos adicionales.
\end{itemize}

De este modo, el auditor dispone de una visión unificada del estado de
seguridad del contrato analizado.

\subsubsection{4.10.7. Resumen del proceso}\label{resumen-del-proceso}

El flujo general implementado por EVMAudit puede resumirse en las
siguientes etapas:

\begin{enumerate}
\def\labelenumi{\arabic{enumi}.}
\tightlist
\item
  Recepción del contrato.
\item
  Configuración del entorno.
\item
  Ejecución de Slither.
\item
  Ejecución de Mythril.
\item
  Normalización.
\item
  Correlación.
\item
  Consulta del catálogo SWC.
\item
  Generación del wrapper.
\item
  Ejecución de Echidna.
\item
  Generación del informe.
\end{enumerate}

La secuencia completa permite combinar distintas técnicas de análisis
dentro de una arquitectura unificada y automatizada.

!TODO: insertar diagrama de secuencia del pipeline completo.

!TODO: insertar diagrama de flujo del proceso.

!TODO: insertar referencia a la Figura X.

!TODO: insertar referencia cruzada a las secciones 4.3 a 4.9.
